Fix distinct groups \(g\neq h\), condition on \(\mathcal F_{b-1}\), and use the uniform perfect-matching representation.  The event that \(g\) is paired with \(h\) has probability \(1/(G-1)\), and its contribution to \(|\mathbb E_0[x_gx_hf_{\eta,\kappa}(W_b)]|\) is therefore at most \(1/(G-1)\).

On the complementary event, \(g\) and \(h\) lie in two different matching pairs.  Conditional on all other orientations, write
\[
 W_b=S+\xi r+\zeta s,
\]
where \(\xi,\zeta\) are the two independent orientations.  From \(|V_{gb}|\leq K_H+2\lambda M\) and \(\sigma_b^2\geq c_\sigma\),
\[
 |r|\vee|s|\leq\frac{2(K_H+2\lambda M)}{\sqrt{Gc_\sigma}}.
\]
The conditional expectation of \(x_gx_hf_{\eta,\kappa}(W_b)\) is, up to a sign, one quarter of the mixed difference
\[
 f(S+r+s)-f(S+r-s)-f(S-r+s)+f(S-r-s).
\]
Twice applying the fundamental theorem of calculus bounds its absolute value by
\(4\|f_{\eta,\kappa}''\|_\infty|rs|\leq C/G\).  Averaging the two matching cases yields
\[
 |\mathbb E_0[x_gx_hf_{\eta,\kappa}(W_b)]|\leq C/G.
\]
Since \(C_b\geq\eta\), division by the soft-kernel normalizer gives
\[
 |\mathbb E_{\mathrm{RR}}(x_{gb}x_{hb}\mid\mathcal F_{b-1})|
 \leq C_\pi/G.
\]

Complement symmetry gives \(\mathbb E_{\mathrm{RR}}x_{gb}=0\).  Thus, for \(s,t\in\{-1,1\}\),
\[
 \Pr(x_{gb}=s,x_{hb}=t\mid\mathcal F_{b-1})
 =\frac14\{1+st\,\mathbb E_{mathrm{RR}}(x_{gb}x_{hb}\mid\mathcal F_{b-1})\}
 =\frac14+O(G^{-1}).
\]
Conditional on the two saturations, the within-group samples are independent across groups.  Multiplication by the appropriate factors \(p_{dq}\) and \(p_{er}\), each at least \(1/4\), proves the uniform product-marginal approximation for every cross-group exposure pair.

For a feasible pair within one group, the first-stage factor is \(1/2\).  The second-stage factor is one of the simple-random-sampling probabilities
\[
 q,\quad 1-q,\quad
 \frac{nq(nq-1)}{n(n-1)},\quad
 \frac{n(1-q)\{n(1-q)-1\}}{n(n-1)},\quad
 \frac{nq\,n(1-q)}{n(n-1)}.
\]
Terms equal to zero are precisely infeasible pairs.  Over \(n\in4\mathbb N\) and \(q\in\{1/4,3/4\}\), every positive displayed term is bounded below by a universal positive constant (the smallest two-treated boundary occurs at \(n=8\)).  This proves the final assertion.